\documentclass{situationist-preprint}

\title{The Possibility of Akrasia}
\msid{ACT-04-a}
\msfamily{ACT-04}
\msdiscipline{Action, Agency and Free Will}
\mstype{Principal manuscript}
\mskeywords{weakness of will; akrasia; better judgment; motivational strength; intentional action; Davidson}

\begin{document}
\maketitle

\begin{abstract}
Can an agent intentionally do one thing while judging, at the time, that it would be better, all things considered, to do another? Socrates denied it, and Davidson's reconstruction of the puzzle shows why the denial is attractive: if an agent who wants one thing more than another does the first when she acts intentionally, and if judging an option better is wanting it more, strict akrasia is impossible. We argue that an agent can intentionally act against her better judgment. The key move is to reject the identification of evaluative ranking with motivational strength. An all-things-considered judgment motivates the agent who makes it, but the strength of the motivation it supplies is not fixed by the judgment's content, and it may be exceeded by the strength of a desire the judgment ranks lower. The resulting action is intentional, since it is done for a reason the agent recognizes, and it is not compelled, since the agent could have resisted. Akrasia is possible, and it is a failure of rationality, not of understanding.
\end{abstract}

\section{The problem}

Weakness of will, or \emph{akrasia}, in the strict sense is free, intentional action contrary to the agent's own judgment, held at the time of action, that some alternative open to her would be better. The woman who judges that she should go to bed, since she must rise early, and who stays up to finish the novel, judging all the while that she should not, is the paradigm. Socrates held that no one does wrong knowingly: an agent who does what is worse believes, at the moment of acting, that it is better, so that apparent akrasia is ignorance \citep[352b--358d]{PlatoProtagoras}. Aristotle allowed that the akratic has knowledge in some sense, but held that it is not in use at the moment of action, overpowered by appetite as in sleep or drunkenness \citep[VII.3]{AristotleNE}. \citet{Hare1952} held that sincerely judging that one ought to do something entails doing it if one can, so that apparent akrasia is either insincere judgment or inability.

\citet{Davidson1970} gave the puzzle its modern form. Two principles seem true. (P1) If an agent wants to do $x$ more than she wants to do $y$, and she believes herself free to do either, she will intentionally do $x$ if she does either intentionally. (P2) If an agent judges that it would be better to do $x$ than $y$, she wants to do $x$ more than $y$. Together they entail that an agent who judges $x$ better and believes herself free to do either does not intentionally do $y$. Davidson's own solution distinguished the all-things-considered judgment, which is conditional on the considerations the agent has weighed, from the unconditional judgment that $x$ is better simpliciter, which is what action expresses. The akratic agent judges that $x$ is better all things considered, but forms the unconditional judgment in favour of $y$. Her failure is a failure to act on the principle of continence, to act on the judgment based on all available relevant reasons.

Davidson's solution makes akrasia possible only by locating the conflict between two judgments, and it has been criticized for making the akratic's unconditional judgment in favour of $y$ an unmotivated mystery \citep{Watson1977,Mele1987}. We defend a simpler solution: P2 is false.

\section{Framework}

\begin{definition}[Better judgment]\label{d:judg}
An agent's \emph{better judgment} at $t$ is her judgment, held at $t$, that a given option open to her is better, all things considered, than the others.
\end{definition}

\begin{definition}[Strict akrasia]\label{d:akr}
An agent acts \emph{akratically} if she intentionally and freely performs an option other than the one her better judgment favours, while holding that judgment.
\end{definition}

\begin{definition}[Motivational strength]\label{d:strength}
The \emph{motivational strength} of an attitude is its causal power to move the agent to act, as measured by what she would do when the attitude conflicts with others.
\end{definition}

\begin{principle}[Pro tanto motivation]\label{p:motiv}
A better judgment motivates the agent who holds it to some degree; its motivational strength is not fixed by its content.
\end{principle}

Principle~\ref{p:motiv} is the internalist thesis that MTE-07 defends, in its weak form: in a rational agent, moral and evaluative judgment necessarily motivates, but motivation is pro tanto and defeasible, so that an agent motivated to some degree by a judgment may be more strongly moved to do something else. The akratic agent is motivated by her judgment, just not enough.

\section{The argument}

\begin{proposition}\label{p:p2}
P2 is false: an agent may judge that $x$ is better than $y$ and want $y$ more than $x$.
\end{proposition}

\begin{proof}
By Definition~\ref{d:strength}, how much an agent wants an option is a matter of the motivational strength of the attitudes that favour it. By Principle~\ref{p:motiv}, a better judgment contributes some strength to the option it favours, but its content, that $x$ is better, does not fix how much. The strength of the desire for $y$ is fixed by other facts: the vividness and proximity of its object, the agent's habits, her fatigue. Nothing in the content of the judgment ensures that its strength exceeds the strength of a desire it ranks lower. There is therefore a possible case in which the agent judges $x$ better and is more strongly moved toward $y$.
\end{proof}

The distinction behind the proposition is the one \citet{Watson1975} drew between \emph{valuing} and \emph{desiring}: the agent's evaluative system ranks options by their worth, and her motivational system ranks them by their pull, and the two systems can diverge. This is not a stipulation made to save the phenomenon. It is confirmed by the familiar fact that the attractiveness of an option varies with its temporal proximity in a way its assessed value does not. The agent who judges in the afternoon that she should go to bed early judges the same at midnight; what has changed is that the novel is open in front of her. This is the pattern of temporal discounting that ACT-13 identifies in procrastination, a form of akrasia.

\begin{theorem}\label{t:main}
Strict akrasia is possible.
\end{theorem}

\begin{proof}
Take the agent of Proposition~\ref{p:p2}, who judges $x$ better than $y$ and wants $y$ more, and who believes herself free to do either. By P1, which we accept, she intentionally does $y$ if she does either. She does it while holding her better judgment, since nothing in the case requires her to revise it. Her action is intentional: she does $y$ for a reason she recognizes, the pleasure of finishing the novel, which is a reason though not a sufficient one. And her action is free, provided she could have resisted the desire, which is the case if her mechanism of action is responsive to reasons in the sense of ACT-01-a: there is a range of situations, for instance those in which the cost of staying up is made more vivid to her, in which the same mechanism would have issued in her going to bed. By Definition~\ref{d:akr}, she acts akratically.
\end{proof}

The proof uses P1 and not P2. We accept P1 because it is close to analytic: what an agent most wants, in the sense of Definition~\ref{d:strength}, is what she is most strongly moved to do.

\section{Objections and replies}

\begin{objection}[Compulsion]
If the desire for $y$ is stronger than the motivation supplied by the judgment, the agent is overpowered, and her action is compelled rather than weak \citep{Watson1977}.
\end{objection}

\begin{reply}
Strength in the sense of Definition~\ref{d:strength} is not irresistibility. A desire is irresistible if no reasons the agent could recognize would lead her to act otherwise; a desire is merely stronger if, in the actual situation, it prevails, though it would not prevail in situations in which she attended to her reasons more closely or exercised the skills of self-control she possesses. \citet{Watson1977} argues that the difference between weakness and compulsion lies in whether the agent has the capacities of self-control that a normal agent has, and we accept this. The akratic agent has the capacity to resist, and fails to exercise it; her failure is culpable for that reason.
\end{reply}

\begin{objection}[Insincerity]
An agent who acts against her judgment does not really hold it; she merely pays it lip service \citep{Hare1952}.
\end{objection}

\begin{reply}
The akratic agent's judgment manifests itself in the ways that sincere judgment does: she regrets her action, criticizes herself, and adjusts her plans for the next occasion. A judgment that has all the marks of sincerity except that it fails to win out on one occasion is sincere. Hare's view requires that judgment be overriding in motivation, which is what Principle~\ref{p:motiv} denies, and the denial is required anyway to make sense of the amoralist, whom MTE-07 counts as possible.
\end{reply}

\begin{objection}[Resolution, not judgment]
\citet{Holton1999} argues that weakness of will, as ordinarily understood, is the unreasonable revision of a resolution, not action against one's better judgment. The puzzle about akrasia concerns a phenomenon that is rare or nonexistent.
\end{objection}

\begin{reply}
Holton is right that ordinary talk of weakness of will often concerns resolutions, and that many cases so called involve no contemporaneous better judgment. But the thesis defended here is not that every case of weakness of will is akrasia; it is that akrasia is possible. That some weak-willed agents abandon a resolution and change their judgment along with it does not show that none acts against a judgment she retains. The woman who stays up while judging that she should not is not revising anything.
\end{reply}

\section{Conclusion}

Davidson's puzzle depends on the identification of evaluative ranking with motivational strength. A better judgment motivates, but the strength of the motivation it supplies is independent of its content, and a desire the judgment ranks lower can be stronger. An agent who acts on that desire while retaining the judgment acts intentionally and, if she could have resisted, freely. Akrasia is possible.

\begin{thebibliography}{99}\small
\bibitem[Aristotle(NE)]{AristotleNE} Aristotle. \emph{Nicomachean Ethics}. Trans. T. Irwin, 2nd ed. Hackett, 1999.
\bibitem[Davidson(1970)]{Davidson1970} Davidson, D. (1970). How is weakness of the will possible? In J. Feinberg (ed.), \emph{Moral Concepts}. Oxford University Press. Reprinted in \emph{Essays on Actions and Events}, Clarendon Press, 1980.
\bibitem[Hare(1952)]{Hare1952} Hare, R.~M. (1952). \emph{The Language of Morals}. Clarendon Press.
\bibitem[Holton(1999)]{Holton1999} Holton, R. (1999). Intention and weakness of will. \emph{The Journal of Philosophy}, 96(5), 241--262.
\bibitem[Mele(1987)]{Mele1987} Mele, A.~R. (1987). \emph{Irrationality: An Essay on Akrasia, Self-Deception, and Self-Control}. Oxford University Press.
\bibitem[Plato(Prot.)]{PlatoProtagoras} Plato. \emph{Protagoras}. Trans. S. Lombardo and K. Bell, in J.~M. Cooper (ed.), \emph{Complete Works}. Hackett, 1997.
\bibitem[Watson(1975)]{Watson1975} Watson, G. (1975). Free agency. \emph{The Journal of Philosophy}, 72(8), 205--220.
\bibitem[Watson(1977)]{Watson1977} Watson, G. (1977). Skepticism about weakness of will. \emph{The Philosophical Review}, 86(3), 316--339.
\end{thebibliography}

\end{document}
